\chapter{Degree two: the theorem of Lefschetz}\label{ch:lefschetz}

In degree two the Hodge conjecture is a theorem, and it holds even with
integral coefficients. The proof takes a page. It is worth giving in
full, because each of its three steps fails in higher degree for a
different reason, and those failures explain why the general case is
hard.

\section{Statement}

\begin{theorem}[Lefschetz]\label{thm:lefschetz11}
Let $X$ be a smooth complex projective variety. Every class
$c\in H^{2}(X,\ZZ)$ whose image in $H^{2}(X,\CC)$ lies in $H^{1,1}(X)$ is
the first Chern class of a line bundle, and therefore the class of a
divisor. In particular $\Hdg^{1}(X)$ is spanned by classes of divisors.
\end{theorem}

\section{Proof}

\subsection*{Step 1: line bundles and the exponential sequence}
On the complex manifold $X$ there is an exact sequence of sheaves
\[
  0\longrightarrow\ZZ\xrightarrow{\;2\pi i\;}\cO_{X}
  \xrightarrow{\;\exp\;}\cO_{X}^{\times}\longrightarrow0 .
\]
Isomorphism classes of holomorphic line bundles are the elements of
$H^{1}(X,\cO_{X}^{\times})$, and the connecting map
$\delta\colon H^{1}(X,\cO_{X}^{\times})\to H^{2}(X,\ZZ)$ of the long exact
sequence is the first Chern class, up to a sign fixed by the
conventions \cite[p.~139]{GH78}. The next terms of the sequence are
\[
  H^{1}(X,\cO_{X}^{\times})\xrightarrow{\;c_{1}\;}H^{2}(X,\ZZ)
  \xrightarrow{\;j\;}H^{2}(X,\cO_{X}),
\]
so an integral class is the Chern class of a holomorphic line bundle
exactly when $j(c)=0$.

\subsection*{Step 2: identifying $j$ with the $(0,2)$ part}
The map $j$ is induced by the inclusion of sheaves $\ZZ\to\cO_{X}$, which
factors as $\ZZ\to\CC\to\cO_{X}$. So $j$ is the composite
\[
  H^{2}(X,\ZZ)\longrightarrow H^{2}(X,\CC)\longrightarrow H^{2}(X,\cO_{X}).
\]
Under the Dolbeault isomorphism $H^{2}(X,\cO_{X})\cong H^{0,2}(X)$, the
second map is the projection of the Hodge decomposition onto its $(0,2)$
summand. To see this, represent a class in $H^{2}(X,\CC)$ by a closed form
$\alpha=\alpha^{2,0}+\alpha^{1,1}+\alpha^{0,2}$ that is harmonic, so that
each component is closed. The map to $H^{2}(X,\cO_{X})$, computed with the
Dolbeault resolution of $\cO_{X}$, keeps only the $(0,2)$ component.

Now let $c$ be an integral class of type $(1,1)$. Its $(0,2)$ component is
zero, so $j(c)=0$, and Step~1 gives a holomorphic line bundle $L$ with
$c_{1}(L)=c$.

The same argument shows something slightly stronger. For a class $c$ with
real image, the $(2,0)$ and $(0,2)$ components are complex conjugate. So
$j(c)=0$ already forces $c$ to be of type $(1,1)$, and the Chern classes of
line bundles are exactly the integral classes of type $(1,1)$.

\subsection*{Step 3: from line bundles to divisors}
Here projectivity enters. By GAGA \cite{Ser56}, the holomorphic line bundle
$L$ on the projective variety $X$ is algebraic. Let $\cO_{X}(1)$ be very
ample. By Serre's theorem, $L(m)=L\otimes\cO_{X}(m)$ is generated by global
sections for $m$ large, and in particular it has a nonzero section, whose
zero scheme is an effective divisor $D_{1}$ with $\cO_{X}(D_{1})\cong L(m)$.
A general hyperplane section $D_{2}$ has $\cO_{X}(D_{2})\cong\cO_{X}(m)$. So
$L\cong\cO_{X}(D_{1}-D_{2})$ and
\[
  c=c_{1}(L)=[D_{1}]-[D_{2}],
\]
using the standard identity $c_{1}(\cO_{X}(D))=[D]$ \cite[p.~141]{GH78}.
This proves the theorem. \qed

\section{Why the proof does not go further}

Each step fails in higher codimension, in a different way.

\subsection*{There is no exponential sequence for cycles of codimension $p\ge2$}
The first step used that line bundles, which are cohomology classes of
$\cO_{X}^{\times}$, are classified by a group that sits in a long exact
sequence with $H^{2}(X,\ZZ)$. The analogue in codimension $p$ is Deligne
cohomology $H^{2p}_{\mathcal D}(X,\ZZ(p))$. It sits in an exact sequence
\[
  0\longrightarrow J^{p}(X)\longrightarrow H^{2p}_{\mathcal D}(X,\ZZ(p))
  \longrightarrow\Hdg^{p}(X)_{\ZZ}\longrightarrow0,
\]
where $J^{p}(X)$ is Griffiths' intermediate Jacobian, and every cycle of
codimension $p$ has a class in it. For $p=1$, Deligne cohomology is
$H^{1}(X,\cO_{X}^{\times})$, the Picard group, and this sequence is the one
of Step~1. For $p\ge2$ the analogous sequence exists and is exact, but it
says nothing about cycles: there is no reason for an element of Deligne
cohomology to come from a cycle. Indeed the image of the cycle map in
$J^{p}(X)$, the image of the Abel--Jacobi map, can be much smaller than
$J^{p}(X)$, as Griffiths showed. Whatever proves the Hodge conjecture
cannot be a cohomological lifting argument of this kind.

\subsection*{Step 2 has no analogue that picks out cycles}
In degree two, being of type $(1,1)$ is equivalent to vanishing of the
$(0,2)$ part, a single condition that the exponential sequence turns into
a geometric object. In degree $2p$ a Hodge class must have vanishing
components of all types $(p+a,p-a)$ with $a\ne0$. Griffiths transversality
then controls how this condition moves in families, but it produces no
object analogous to a line bundle.

\subsection*{Step 3 has no analogue at all}
Even on a compact K\"ahler manifold, Steps~1 and~2 go through, and every
integral $(1,1)$ class is the Chern class of a holomorphic line bundle. It
is only Step~3 that uses projectivity, and on a non-algebraic complex torus
the line bundle produced in Step~1 may have no meromorphic sections at all.
In higher codimension there is not even a line bundle to start from. The
nearest substitute is a coherent sheaf, or a vector bundle, with prescribed
Chern character. On a very general complex torus of Weil type, Voisin
showed that no such sheaf exists \cite{Voi02b}. On an abelian variety we do
not know.

These three failures explain the shape of what follows. Since there is no
general mechanism for producing cycles, the known results for abelian
varieties all start from divisors, the only cycles we can always produce,
and move them around with correspondences. Chapter~\ref{ch:abelian} shows
how far that takes us.
